\Chapter{20}

\Verse{1}
{
    \zA{话说宝玉在黛玉房中说“耗子精”，宝钗撞来，讽刺宝玉元宵不知“绿蜡”之 典，三人正在房中互相取笑。}
}
{
    \eA{But to continue. Pao-yü was in Tai yü's apartments relating about the
        rat-elves, when Pao-ch'ai entered unannounced, and began to gibe Pao-yü,
        with trenchant irony: how that on the fifteenth of the first moon, he had
        shown ignorance of the allusion to the green wax; and the three of them then
        indulged in that room in mutual poignant satire,}
}
{
}

\Verse{2}
{
    \zA{那宝玉恐黛玉饭后贪眠，一时存了食，或夜间走了困， 身体不好；幸而宝钗走来，大家谈笑，那黛玉方不欲睡，自己才放了心。}
}
{
    \eA{for the sake of fun. Pao-yü had been giving way to solicitude lest Tai-yü
        should, by being bent upon napping soon after her meal, be shortly getting
        an indigestion, or lest sleep should, at night, be completely dispelled, as
        neither of these things were conducive to the preservation of good health,
        when luckily Pao-ch'ai walked in,}
}
{
}

\Verse{3}
{
    \zA{忽听他房 中嚷起来，大家侧耳听了一听，黛玉先笑道：“这是你妈妈和袭人叫唤呢。}
}
{
    \eA{and they chatted and laughed together; and when Lin Tai-yü at length lost
        all inclination to dose, he himself then felt composed in his mind. But
        suddenly they heard clamouring begin in his room, and after they had all
        lent an ear and listened, Lin Tai-yü was the first to smile and make a
        remark. "It's your nurse having a row with Hsi Jen!"}
}
{
}

\Verse{4}
{
    \zA{那袭人 待他也罢了，你妈妈再要认真排揎他，可见老背晦了。”}
}
{
    \eA{she said. "Hsi Jen treats her well enough, but that nurse of yours would
        also like to keep her well under her thumb; she's indeed an old dotard;" and
        Pao-yü was anxious to go over at once,}
}
{
}

\Verse{5}
{
    \zA{宝玉忙欲赶过去，宝钗一 把拉住道：“你别和你妈妈吵才是呢!他是老糊涂了，倒要让他一步儿的是。”}
}
{
    \eA{but Pao-ch'ai laid hold of him and kept him back, suggesting: "It's as well
        that you shouldn't wrangle with your nurse, for she's quite stupid from old
        age; and it's but fair, on the contrary, that you should bear with her a
        little." "I know all about that!"}
}
{
}

\Verse{6}
{
    \zA{宝 玉道：“我知道了。”}
}
{
    \eA{Pao-yü rejoined. But having concluded this remark, he walked into his room,}
}
{
}

\Verse{7}
{
    \zA{说毕走来。}
}
{
    \eA{where he discovered nurse Li,}
}
{
}

\Verse{8}
{
    \zA{只见李嬷嬷拄着拐杖，在当地骂袭人：“忘了本的小娼妇儿!我抬举起你来， 这会子我来了，你大模厮样儿的躺在炕上，见了我也不理一理儿。}
}
{
    \eA{leaning on her staff, standing in the centre of the floor, abusing Hsi Jen,
        saying: "You young wench! how utterly unmindful you are of your origin! It's
        I who've raised you up, and yet, when I came just now, you put on high airs
        and mighty side,}
}
{
}

\Verse{9}
{
    \zA{一心只想妆狐媚 子哄宝玉，哄的宝玉不理我，只听你的话。}
}
{
    \eA{and remained reclining on the stove-couch! You saw me well enough, but you
        paid not the least heed to me! Your whole heart is set upon acting like a
        wily enchantress to befool Pao-yü;}
}
{
}

\Verse{10}
{
    \zA{你不过是几两银子买了来的小丫头子罢 咧，这屋里你就作起耗来了!好不好的，拉出去配一个小子，看你还妖精似的哄人 不哄！”}
}
{
    \eA{and you so impose upon Pao-yü that he doesn't notice me, but merely lends an
        ear to what you people have to say! You're no more than a low girl bought
        for a few taels and brought in here; and will it ever do that you should be
        up to your mischievous tricks in this room? But whether you like it or not,
        I'll drag you out from this, and give you to some mean fellow,}
}
{
}

\Verse{11}
{
    \zA{袭人先只道李嬷嬷不过因他躺着生气，少不得分辩说：“病了，才出汗， 蒙着头，原没看见你老人家。”}
}
{
    \eA{and we'll see whether you will still behave like a very imp, and cajole
        people or not?" Hsi Jen was, at first, under the simple impression that the
        nurse was wrath for no other reason than because she remained lying down,
        and she felt constrained to explain that "she was unwell, that she had just
        succeeded in perspiring,}
}
{
}

\Verse{12}
{
    \zA{后来听见他说“哄宝玉”，又说“配小子”，由不 得又羞又委屈，禁不住哭起来了。}
}
{
    \eA{and that having had her head covered, she hadn't really perceived the old
        lady;" but when she came subsequently to hear her mention that she imposed
        upon Pao-yü, and also go so far as to add that she would be given to some
        mean fellow,}
}
{
}

\Verse{13}
{
    \zA{宝玉虽听了这些话，也不好怎样，少不得替他分 辩，说“病了，吃药”，又说：“你不信，只问别的丫头。”}
}
{
    \eA{she unavoidably experienced both a sense of shame and injury, and found it
        impossible to restrain herself from beginning to cry. Pao-yü had, it is
        true, caught all that had been said, but unable with any propriety to take
        notice of it, he thought it his duty to explain matters for her.}
}
{
}

\Verse{14}
{
    \zA{李嬷嬷听了这话，越 发气起来了，说道：“你只护着那起狐狸，那里还认得我了呢？}
}
{
    \eA{"She's ill," he observed, "and is taking medicines; and if you don't believe
        it," he went on, "well then ask the rest of the servant-girls." Nurse Li at
        these words flew into a more violent dudgeon. "Your sole delight is to
        screen that lot of sly foxes!"}
}
{
}

\Verse{15}
{
    \zA{叫我问谁去？}
}
{
    \eA{she remarked, "and do you pay any notice to me?}
}
{
}

\Verse{16}
{
    \zA{谁不帮 着你呢？}
}
{
    \eA{No, none at all! and whom would you like me to go and ask;}
}
{
}

\Verse{17}
{
    \zA{谁不是袭人拿下马来的？}
}
{
    \eA{who's it that doesn't back you? and who hasn't been dismounted from her
        horse by Hsi Jen?}
}
{
}

\Verse{18}
{
    \zA{我都知道那些事!我只和你到老太太、太太跟前去 讲讲：把你奶了这么大，到如今吃不着奶了，把我扔在一边儿，逞着丫头们要我的 强！”}
}
{
    \eA{I know all about it; but I'll go with you and explain all these matters to
        our old mistress and my lady; for I've nursed you till I've brought you to
        this age, and now that you don't feed on milk, you thrust me on one side,
        and avail yourself of the servant-girls, in your wish to browbeat me."}
}
{
}

\Verse{19}
{
    \zA{一面说，一面哭。}
}
{
    \eA{As she uttered this remark, she too gave way to tears,}
}
{
}

\Verse{20}
{
    \zA{彼时黛玉宝钗等也过来劝道：“妈妈，你老人家担待他们 些就完了。”}
}
{
    \eA{but by this time, Tai-yü and Pao-ch'ai had also come over, and they set to
        work to reassure her. "You, old lady," they urged, "should bear with them a
        little, and everything will be right!"}
}
{
}

\Verse{21}
{
    \zA{李嬷嬷见他二人来了，便诉委屈，将当日吃茶，茜雪出去，和昨日酥 酪等事，唠唠叨叨说个不了。}
}
{
    \eA{And when nurse Li saw these two arrive, she hastened to lay bare her
        grievances to them; and taking up the question of the dismissal in days gone
        by, of Hsi Hsüeh, for having drunk some tea, of the cream eaten on the
        previous day, and other similar matters,}
}
{
}

\Verse{22}
{
    \zA{可巧凤姐正在上房算了输赢帐，听见后面一片声嚷，便知是李嬷嬷老病发了， 又值他今儿输了钱，迁怒于人，排揎宝玉的丫头。}
}
{
    \eA{she spun a long, interminable yarn. By a strange coincidence lady Feng was
        at this moment in the upper rooms, where she had been making up the account
        of losses and winnings, and upon hearing at the back a continuous sound of
        shouting and bustling, she readily concluded that nurse Li's old complaint
        was breaking forth, and that she was finding fault with Pao-yü's servants.}
}
{
}

\Verse{23}
{
    \zA{便连忙赶过来拉了李嬷嬷，笑道： “妈妈别生气。}
}
{
    \eA{But she had, as luck would have it, lost money in gambling on this occasion,
        so that she was ready to visit her resentment upon others.}
}
{
}

\Verse{24}
{
    \zA{大节下，老太太刚喜欢了一日。}
}
{
    \eA{With hurried step, she forthwith came over, and laying hold of nurse Li,
        "Nurse," she said smiling,}
}
{
}

\Verse{25}
{
    \zA{你是个老人家，别人吵，你还要管 他们才是；难道你倒不知规矩，在这里嚷起来，叫老太太生气不成？}
}
{
    \eA{"don't lose your temper, on a great festival like this, and after our
        venerable lady has just gone through a day in excellent spirits! You're an
        old dame, and should, when others get up a row, still do what is right and
        keep them in proper order; and aren't you, instead of that,}
}
{
}

\Verse{26}
{
    \zA{你说谁不好， 我替你打他。}
}
{
    \eA{aware what good manners imply, that you will start vociferating in this
        place,}
}
{
}

\Verse{27}
{
    \zA{我屋里烧的滚热的野鸡，快跟了我喝酒去罢。”}
}
{
    \eA{and make our dowager lady full of displeasure? Tell me who's not good, and
        I'll beat her for you; but be quick and come along with me over to my
        quarters,}
}
{
}

\Verse{28}
{
    \zA{一面说，一面拉着走， 又叫：“丰儿，替你李奶奶拿着拐棍子、擦眼泪的绢子。”}
}
{
    \eA{where a pheasant which they have roasted is scalding hot, and let us go and
        have a glass of wine!" And as she spoke, she dragged her along and went on
        her way. "Feng Erh," she also called, "hold the staff for your old lady Li,}
}
{
}

\Verse{29}
{
    \zA{那李嬷嬷脚不沾地跟了 凤姐儿走了，一面还说：“我也不要这老命了，索性今儿没了规矩，闹一场子，讨 了没脸，强似受那些娼妇的气！”}
}
{
    \eA{and the handkerchief to wipe her tears with!" While nurse Li walked along
        with lady Feng, her feet scarcely touched the ground, as she kept on saying:
        "I don't really attach any value to this decrepid existence of mine! and I
        had rather disregard good manners, have a row and lose face, as it's better,
        it seems to me,}
}
{
}

\Verse{30}
{
    \zA{后面宝钗黛玉见凤姐儿这般，都拍手笑道：“亏 他这一阵风来，把个老婆子撮了去了。”}
}
{
    \eA{than to put up with the temper of that wench!" Behind followed Pao-ch'ai and
        Tai-yü, and at the sight of the way in which lady Feng dealt with her, they
        both clapped their hands, and exclaimed, laughing, "What piece of luck that
        this gust of wind has come,}
}
{
}

\Verse{31}
{
    \zA{宝玉点头叹道：“这又不知是那里的帐，只拣软的欺负!又不知是那个姑娘得 罪了，上在他帐上了。”}
}
{
    \eA{and dragged away this old matron!" while Pao-yü nodded his head to and fro
        and soliloquised with a sigh: "One can neither know whence originates this
        score; for she will choose the weak one to maltreat; nor can one see what
        girl has given her offence that she has come to be put in her black books!"}
}
{
}

\Verse{32}
{
    \zA{一句未完，晴雯在旁说道：“谁又没疯了，得罪他做什么？}
}
{
    \eA{Scarcely had he ended this remark, before Ch'ing Wen, who stood by, put in
        her word. "Who's gone mad again?" she interposed, "and what good would come
        by hurting her feelings?}
}
{
}

\Verse{33}
{
    \zA{既得罪了他，就有本事承任，犯不着带累别人！”}
}
{
    \eA{But did even any one happen to hurt her, she would have pluck enough to bear
        the brunt, and wouldn't act so improperly as to involve others!"}
}
{
}

\Verse{34}
{
    \zA{袭人一面哭，一面拉着宝玉道： “为我得罪了一个老奶奶，你这会子又为我得罪这些人，这还不够我受的，还只是 拉扯人！”}
}
{
    \eA{Hsi Jen wept, and as she, did so, she drew Pao-yü towards her: "All through
        my having aggrieved an old nurse," she urged, "you've now again given
        umbrage, entirely on my account, to this crowd of people; and isn't this
        still enough for me to bear but must you also go and drag in third parties?"}
}
{
}

\Verse{35}
{
    \zA{宝玉见他这般病势，又添了这些烦恼，连忙忍气吞声，安慰他仍旧睡下 出汗。}
}
{
    \eA{When Pao-yü realised that to this sickness of hers, had also been superadded
        all these annoyances, he promptly stifled his resentment, suppressed his
        voice and consoled her so far as to induce her to lie down again to
        perspire.}
}
{
}

\Verse{36}
{
    \zA{又见他汤烧火热，自己守着他，歪在旁边，劝他只养病，别想那些没要紧的 事。}
}
{
    \eA{And when he further noticed how scalding like soup and burning like fire she
        was, he himself watched by her, and reclining by her side, he tried to cheer
        her, saying: "All you must do is to take good care of your ailment;}
}
{
}

\Verse{37}
{
    \zA{袭人冷笑道：“要为这些事生气，这屋里一刻还住得了？}
}
{
    \eA{and don't give your mind to those trifling matters, and get angry." "Were
        I," Hsi Jen smiled sardonically, "to lose my temper over such concerns,}
}
{
}

\Verse{38}
{
    \zA{但只是天长日久，尽 着这么闹，可叫人怎么过呢!你只顾一时为我得罪了人，他们都记在心里，遇着坎 儿，说的好说不好听的，大家什么意思呢？”}
}
{
    \eA{would I be able to stand one moment longer in this room? The only thing is
        that if she goes on, day after day, doing nothing else than clamour in this
        manner, how can she let people get along? But you rashly go and hurt
        people's feelings for our sakes; but they'll bear it in mind, and when they
        find an opportunity, they'll come out with what's easy enough to say,}
}
{
}

\Verse{39}
{
    \zA{一面说，一面禁不住流泪，又怕宝玉 烦恼，只得又勉强忍着。}
}
{
    \eA{but what's not pleasant to hear, and how will we all feel then?" While her
        mouth gave utterance to these words, she could not stop her tears from
        running; but fearful, on the other hand,}
}
{
}

\Verse{40}
{
    \zA{一时杂使的老婆子端了二和药来，宝玉见他才有点汗儿， 便不叫他起来，自己端着给他就枕上吃了，即令小丫鬟们铺炕。}
}
{
    \eA{lest Pao-yü should be annoyed, she felt compelled to again strain every
        nerve to repress them. But in a short while, the old matrons employed for
        all sorts of duties, brought in some mixture of two drugs; and, as Pao-yü
        noticed that she was just on the point of perspiring, he did not allow her
        to get up, but readily taking it up to her,}
}
{
}

\Verse{41}
{
    \zA{袭人道：“你吃饭 不吃饭，到底老太太、太太跟前坐一会子，和姑娘们玩一会子，再回来。}
}
{
    \eA{she immediately swallowed it, with her head still on her pillow; whereupon
        he gave speedy directions to the young servant-maids to lay her stove-couch
        in order. "Whether you mean to have anything to eat or not," Hsi Jen
        advised, "you should after all sit for a time with our old mistress and our
        lady,}
}
{
}

\Verse{42}
{
    \zA{我就静静 的躺一躺也好啊。”}
}
{
    \eA{and have a romp with the young ladies; after which you can come back again;}
}
{
}

\Verse{43}
{
    \zA{宝玉听说，只得依他，看着他去了簪环躺下，才去上屋里跟着 贾母吃饭。}
}
{
    \eA{while I, by quietly keeping lying down, will also feel the better." When
        Pao-yü heard this suggestion, he had no help but to accede, and, after she
        had divested herself of her hair-pins and earrings,}
}
{
}

\Verse{44}
{
    \zA{饭毕，贾母犹欲和那几个老管家的嬷嬷斗牌。}
}
{
    \eA{and he saw her lie down, he betook himself into the drawing-rooms, where he
        had his repast with old lady Chia. But the meal over, her ladyship felt
        still disposed to play at cards with the nurses,}
}
{
}

\Verse{45}
{
    \zA{宝玉惦记袭人，便回至房中。}
}
{
    \eA{who had looked after the household for many years; and Pao-yü, bethinking
        himself of Hsi Jen,}
}
{
}

\Verse{46}
{
    \zA{见 袭人朦胧睡去，自己要睡，天气尚早。}
}
{
    \eA{hastened to return to his apartments; where seeing that Hsi Jen was drowsily
        falling asleep, he himself would have wished to go to bed,}
}
{
}

\Verse{47}
{
    \zA{彼时晴雯、绮霞、秋纹、碧痕都寻热闹，找 鸳鸯、琥珀等耍戏去了。}
}
{
    \eA{but the hour was yet early. And as about this time Ch'ing Wen, I Hsia, Ch'in
        Wen, Pi Hen had all, in their desire of getting some excitement, started in
        search of Yüan Yang,}
}
{
}

\Verse{48}
{
    \zA{见麝月一人在外间屋里灯下抹骨牌。}
}
{
    \eA{Hu Po and their companions, to have a romp with them, and he espied She Yüeh
        alone in the outer room,}
}
{
}

\Verse{49}
{
    \zA{宝玉笑道：“你怎么 不和他们去？”}
}
{
    \eA{having a game of dominoes by lamp-light, Pao-yü inquired full of smiles:
        "How is it you don't go with them?"}
}
{
}

\Verse{50}
{
    \zA{麝月道：“没有钱。”}
}
{
    \eA{"I've no money," She Yüeh replied. "Under the bed," continued Pao-yü,}
}
{
}

\Verse{51}
{
    \zA{宝玉道：“床底下堆着钱，还不够你输的？”}
}
{
    \eA{"is heaped up all that money, and isn't it enough yet for you to lose from?"
        "Had we all gone to play,"}
}
{
}

\Verse{52}
{
    \zA{麝月道：“都乐去了，这屋子交给谁呢？}
}
{
    \eA{She Yüeh added, "to whom would the charge of this apartment have been handed
        over? That other one is sick again,}
}
{
}

\Verse{53}
{
    \zA{那一个又病了，满屋里上头是灯，下头是 火，那些老婆子们都老天拔地伏侍了一天，也该叫他们歇歇儿了。}
}
{
    \eA{and the whole room is above, one mass of lamps, and below, full of fire; and
        all those old matrons, ancient as the heavens, should, after all their
        exertions in waiting upon you from morning to night, be also allowed some
        rest;}
}
{
}

\Verse{54}
{
    \zA{小丫头们也伏侍 了一天，这会子还不叫玩玩儿去吗？}
}
{
    \eA{while the young servant girls, on the other hand, have likewise been on duty
        the whole day long, and shouldn't they even at this hour be left to go and
        have some distraction?}
}
{
}

\Verse{55}
{
    \zA{所以我在这里看着。”}
}
{
    \eA{and that's why I am in here on watch." When Pao-yü heard these words,}
}
{
}

\Verse{56}
{
    \zA{宝玉听了这话，公然又 是一个袭人了。}
}
{
    \eA{which demonstrated distinctly that she was another Hsi Jen, he consequently
        put on a smile and remarked:}
}
{
}

\Verse{57}
{
    \zA{因笑道：“我在这里坐着，你放心去罢。”}
}
{
    \eA{"I'll sit in here, so you had better set your mind at ease and go!" "Since
        you remain in here,}
}
{
}

\Verse{58}
{
    \zA{麝月道：“你既在这里， 越发不用去了。}
}
{
    \eA{there's less need for me to go," resumed She Yüeh, "for we two can chat and
        play and laugh; and won't that be nice?"}
}
{
}

\Verse{59}
{
    \zA{咱们两个说话儿不好？”}
}
{
    \eA{"What can we two do? it will be awfully dull! but never mind," Pao-yü
        rejoined;}
}
{
}

\Verse{60}
{
    \zA{宝玉道：“咱们两个做什么呢？}
}
{
    \eA{"this morning you said that your head itched, and now that you have nothing
        to do,}
}
{
}

\Verse{61}
{
    \zA{怪没意思 的。}
}
{
    \eA{I may as well comb it for you."}
}
{
}

\Verse{62}
{
    \zA{也罢了，早起你说头上痒痒，这会子没什么事，我替你篦头罢。”}
}
{
    \eA{"Yes! do so!" readily assented She Yüeh, upon catching what he suggested;
        and while still speaking, she brought over the dressing-case containing a
        set of small drawers and looking-glass,}
}
{
}

\Verse{63}
{
    \zA{麝月听了道： “使得。”}
}
{
    \eA{and taking off her ornaments, she dishevelled her hair;}
}
{
}

\Verse{64}
{
    \zA{说着，将文具镜匣搬来，卸去钗？}
}
{
    \eA{whereupon Pao-yü picked up the fine comb and passed it repeatedly through
        her hair;}
}
{
}

\Verse{65}
{
    \zA{，打开头发，宝玉拿了篦子替他篦。}
}
{
    \eA{but he had only combed it three or five times, when he perceived Ch'ing Wen
        hurriedly walk in to fetch some money.}
}
{
}

\Verse{66}
{
    \zA{只篦了三五下儿，见晴雯忙忙走进来取钱，一见他两个，便冷笑道：“哦!交 杯盏儿还没吃，就上了头了！”}
}
{
    \eA{As soon as she caught sight of them both: "You haven't as yet drunk from the
        marriage cup," she said with a smile full of irony, "and have you already
        put up your hair?" "Now that you've come, let me also comb yours for you,"
        Pao-yü continued. "I'm not blessed with such excessive good fortune!"}
}
{
}

\Verse{67}
{
    \zA{宝玉笑道：“你来，我也替你篦篦。”}
}
{
    \eA{Ch'ing Wen retorted, and as she uttered these words, she took the money, and
        forthwith dashing the portiere after her,}
}
{
}

\Verse{68}
{
    \zA{晴雯道：“我 没这么大造化。”}
}
{
    \eA{she quitted the room. Pao-yü stood at the back of She Yüeh, and She Yüeh sat
        opposite the glass,}
}
{
}

\Verse{69}
{
    \zA{说着，拿了钱，摔了帘子，就出去了。}
}
{
    \eA{so that the two of them faced each other in it, and Pao-yü readily observed
        as he gazed in the glass,}
}
{
}

\Verse{70}
{
    \zA{宝玉在麝月身后，麝月对 镜，二人在镜内相视而笑。}
}
{
    \eA{"In the whole number of rooms she's the only one who has a glib tongue!" She
        Yüeh at these words hastily waved her hand towards the inside of the glass,}
}
{
}

\Verse{71}
{
    \zA{宝玉笑着道：“满屋里就只是他磨牙。”}
}
{
    \eA{and Pao-yü understood the hint; and suddenly a sound of "hu" was heard from
        the portiere,}
}
{
}
\ChapterImage{img/chapters/067第20-21回 林黛玉俏語谑嬌音 俊襲人嬌箴嗔寶玉 俏平兄軟語庇賈璉.jpg}


\Verse{72}
{
    \zA{麝月听说，忙 向镜中摆手儿。}
}
{
    \eA{and Ch'ing Wen ran in once again. "How have I got a glib tongue?"}
}
{
}

\Verse{73}
{
    \zA{宝玉会意，忽听“唿”一声帘子响，晴雯又跑进来问道：“我怎么 磨牙了？}
}
{
    \eA{she inquired; "it would be well for us to explain ourselves." "Go after your
        business, and have done," She Yüeh interposed laughingly; "what's the use of
        your coming and asking questions of people?"}
}
{
}

\Verse{74}
{
    \zA{咱们倒得说说！”}
}
{
    \eA{"Will you also screen him?" Ch'ing Wen smiled significantly;}
}
{
}

\Verse{75}
{
    \zA{麝月笑道：“你去你的罢，又来拌嘴儿了。”}
}
{
    \eA{"I know all about your secret doings, but wait until I've got back my
        capital, and we'll then talk matters over!"}
}
{
}

\Verse{76}
{
    \zA{晴雯也笑 道：“你又护着他了!你们瞒神弄鬼的，打量我都不知道呢!等我捞回本儿来再说。”}
}
{
    \eA{With this remark still on her lips, she straightway quitted the room, and
        during this while, Pao-yü having finished combing her hair, asked She Yüeh
        to quietly wait upon him, while he went to sleep, as he would not like to
        disturb Hsi Jen.}
}
{
}

\Verse{77}
{
    \zA{说着，一径去了。}
}
{
    \eA{Of the whole night there is nothing to record.}
}
{
}

\Verse{78}
{
    \zA{这里宝玉通了头，命麝月悄悄的伏侍他睡下，不肯惊动袭人。}
}
{
    \eA{But the next day, when he got up at early dawn, Hsi Jen had already
        perspired, during the night, so that she felt considerably lighter and
        better;}
}
{
}

\Verse{79}
{
    \zA{一 宿无话。}
}
{
    \eA{but limiting her diet to a little rice soup,}
}
{
}

\Verse{80}
{
    \zA{次日清晨，袭人已是夜间出了汗，觉得轻松了些，只吃些米汤静养。}
}
{
    \eA{she remained quiet and nursed herself, and Pao-yü was so relieved in mind
        that he came, after his meal, over on this side to his aunt Hsüeh's on a
        saunter. The season was the course of the first moon,}
}
{
}

\Verse{81}
{
    \zA{宝玉才放 了心，因饭后走到薛姨妈这边来闲逛。}
}
{
    \eA{and the school was shut up for the new year holidays; while in the inner
        chambers the girls had put by their needlework, and were all having a time
        of leisure,}
}
{
}

\Verse{82}
{
    \zA{彼时正月内学房中放年学，闺阁中忌针黹，都是闲时，因贾环也过来玩。}
}
{
    \eA{and hence it was that when Chia Huan too came over in search of distraction,
        he discovered Pao-ch'ai, Hsiang Ling, Ying Erh, the three of them, in the
        act of recreating themselves by playing at chess.}
}
{
}

\Verse{83}
{
    \zA{正遇 见宝钗、香菱、莺儿三个赶围棋作耍，贾环见了也要玩。}
}
{
    \eA{Chia Huan, at the sight of them, also wished to join in their games; and
        Pao-ch'ai, who had always looked upon him with, in fact, the same eye as she
        did Pao-yü,}
}
{
}

\Verse{84}
{
    \zA{宝钗素日看他也如宝玉， 并没他意，今儿听他要玩，让他上来，坐在一处玩。}
}
{
    \eA{and with no different sentiment of any kind, pressed him to come up, upon
        hearing that he was on this occasion desirous to play; and, when he had
        seated himself together with them, they began to gamble, staking each time a
        pile of ten cash.}
}
{
}

\Verse{85}
{
    \zA{一注十个钱。}
}
{
    \eA{The first time, he was the winner,}
}
{
}

\Verse{86}
{
    \zA{头一回，自己赢 了，心中十分喜欢。}
}
{
    \eA{and he felt supremely elated at heart, but as it happened that he
        subsequently lost in several consecutive games he soon became a prey to
        considerable distress.}
}
{
}

\Verse{87}
{
    \zA{谁知后来接连输了几盘，就有些着急。}
}
{
    \eA{But in due course came the game in which it was his turn to cast the dice,
        and, if in throwing, he got seven spots,}
}
{
}

\Verse{88}
{
    \zA{赶着这盘正该自己掷骰 子，若掷个七点便赢了，若掷个六点也该赢，掷个三点就输了。}
}
{
    \eA{he stood to win, but he was likewise bound to be a winner were he to turn up
        six; and when Ying Erh had turned up three spots and lost, he consequently
        took up the dice, and dashing them with spite, one of them settled at five;}
}
{
}

\Verse{89}
{
    \zA{因拿起骰子来狠命 一掷，一个坐定了二，那一个乱转。}
}
{
    \eA{and, as the other reeled wildly about, Ying Erh clapped her hands, and kept
        on shouting, "one spot;" while Chia Huan at once gazed with fixed eye and
        cried at random:}
}
{
}

\Verse{90}
{
    \zA{莺儿拍着手儿叫“么！”}
}
{
    \eA{"It's six, it's seven, it's eight!" But the dice, as it happened,}
}
{
}

\Verse{91}
{
    \zA{贾环便瞪着眼，“六！”}
}
{
    \eA{turned up at one spot, and Chia Huan was so exasperated that putting out his
        hand,}
}
{
}

\Verse{92}
{
    \zA{“七！”}
}
{
    \eA{he speedily made a snatch at the dice,}
}
{
}

\Verse{93}
{
    \zA{“八！”}
}
{
    \eA{and eventually was about to lay hold of the money,}
}
{
}

\Verse{94}
{
    \zA{混叫。}
}
{
    \eA{arguing that it was six spot.}
}
{
}

\Verse{95}
{
    \zA{那骰子偏生转出么来。}
}
{
    \eA{But Ying Erh expostulated, "It was distinctly an ace,"}
}
{
}

\Verse{96}
{
    \zA{贾环急了，伸手便抓起骰子来，就 要拿钱，说是个四点。}
}
{
    \eA{she said. And as Pao-ch'ai noticed how distressed Chia Huan was, she
        forthwith cast a glance at Ying Erh and observed: "The older you get, the
        less manners you have!}
}
{
}

\Verse{97}
{
    \zA{莺儿便说：“明明是个么！”}
}
{
    \eA{Is it likely that gentlemen will cheat you? and don't you yet put down the
        money?"}
}
{
}

\Verse{98}
{
    \zA{宝钗见贾环急了，便瞅了莺儿 一眼，说道：“越大越没规矩!难道爷们还赖你？}
}
{
    \eA{Ying Erh felt her whole heart much aggrieved, but as she heard Pao-ch'ai
        make these remarks, she did not presume to utter a sound, and as she was
        under the necessity of laying down the cash,}
}
{
}

\Verse{99}
{
    \zA{还不放下钱来呢。”}
}
{
    \eA{she muttered to herself: "This one calls himself a gentleman,}
}
{
}

\Verse{100}
{
    \zA{莺儿满心委屈， 见姑娘说，不敢出声，只得放下钱来，口内嘟囔说：“一个做爷的，还赖我们这几
        个钱，连我也瞧不起!前儿和宝二爷玩，他输了那些也没着急，下剩的钱还是几个 小丫头子们一抢，他一笑就罢了。”}
}
{
    \eA{and yet cheats us of these few cash, for which I myself even have no eye!
        The other day when I played with Mr. Pao-yü, he lost ever so many, and yet
        he did not distress himself! and what remained of the cash were besides
        snatched away by a few servant-girls,}
}
{
}

\Verse{101}
{
    \zA{宝钗不等说完，连忙喝住了。}
}
{
    \eA{but all he did was to smile, that's all!" Pao-ch'ai did not allow her time
        to complete what she had to say,}
}
{
}

\Verse{102}
{
    \zA{贾环道：“我拿什么比宝玉？}
}
{
    \eA{but there and then called her to account and made her desist; whereupon Chia
        Huan exclaimed:}
}
{
}

\Verse{103}
{
    \zA{你们怕他，都和他 好，都欺负我不是太太养的！”}
}
{
    \eA{"How can I compare with Pao-yü; you all fear him, and keep on good terms
        with him, while you all look down upon me for not being the child of my
        lady."}
}
{
}

\Verse{104}
{
    \zA{说着便哭。}
}
{
    \eA{And as he uttered these words,}
}
{
}

\Verse{105}
{
    \zA{宝钗忙劝他：“好兄弟，快别说这话， 人家笑话。”}
}
{
    \eA{he at once gave way to tears. "My dear cousin," Pao-ch'ai hastened to advise
        him, "leave off at once language of this kind,}
}
{
}

\Verse{106}
{
    \zA{又骂莺儿。}
}
{
    \eA{for people will laugh at you;"}
}
{
}

\Verse{107}
{
    \zA{正值宝玉走来，见了这般景况，问：“是怎么了？”}
}
{
    \eA{and then went on to scold Ying Erh, when Pao-yü just happened to come in.
        Perceiving him in this plight, "What is the matter?" he asked;}
}
{
}

\Verse{108}
{
    \zA{贾环 不敢则声。}
}
{
    \eA{but Chia Huan had not the courage to say anything.}
}
{
}

\Verse{109}
{
    \zA{宝钗素知他家规矩，凡做兄弟的怕哥哥。}
}
{
    \eA{Pao-ch'ai was well aware of the custom, which prevailed in their family,
        that younger brothers lived in respect of the elder brothers,}
}
{
}

\Verse{110}
{
    \zA{却不知那宝玉是不要人怕他的， 他想着：“兄弟们一并都有父母教训，何必我多事，反生疏了。}
}
{
    \eA{but she was not however cognisant of the fact that Pao-yü would not that any
        one should entertain any fear of him. His idea being that elder as well as
        younger brothers had, all alike, father and mother to admonish them, and
        that there was no need for any of that officiousness,}
}
{
}

\Verse{111}
{
    \zA{况且我是正出，他 是庶出，饶这样看待，还有人背后谈论，还禁得辖治了他？”}
}
{
    \eA{which, instead of doing good gave, on the contrary, rise to estrangement.
        "Besides," (he reasoned,) "I'm the offspring of the primary wife, while he's
        the son of the secondary wife,}
}
{
}

\Verse{112}
{
    \zA{更有个呆意思存在心 里。}
}
{
    \eA{and, if by treating him as leniently as I have done, there are still those
        to talk about me,}
}
{
}

\Verse{113}
{
    \zA{你道是何呆意？}
}
{
    \eA{behind my back, how could I exercise any control over him?"}
}
{
}

\Verse{114}
{
    \zA{因他自幼姐妹丛中长大，亲姊妹有元春探春，叔伯的有迎春惜 春，亲戚中又有湘云黛玉宝钗等人，他便料定天地间灵淑之气只钟于女子，男儿们 不过是些渣滓浊沫而已。}
}
{
    \eA{But besides these, there were other still more foolish notions, which he
        fostered in his mind; but what foolish notions they were can you, reader,
        guess? As a result of his growing up, from his early youth, among a crowd of
        girls, of whom, in the way of sister, there was Yüan Ch'un, of cousins,}
}
{
}

\Verse{115}
{
    \zA{因此把一切男子都看成浊物，可有可无。}
}
{
    \eA{from his paternal uncle's side, there were Ying Ch'un, and Hsi Ch'un, and of
        relatives also there were Shih Hsiang-yün,}
}
{
}

\Verse{116}
{
    \zA{只是父亲、伯叔、 兄弟之伦，因是圣人遗训，不敢违忤，所以弟兄间亦不过尽其大概就罢了，并不想 自己是男子，须要为子弟之表率。}
}
{
    \eA{Lin Tai-yü, Hsüeh Pao-ch'ai and the rest, he, in due course, resolved in his
        mind that the divine and unsullied virtue of Heaven and earth was only
        implanted in womankind, and that men were no more than feculent dregs and
        foul dirt. And for this reason it was that men were without discrimination,
        considered by him as so many filthy objects, which might or might not exist;}
}
{
}

\Verse{117}
{
    \zA{是以贾环等都不甚怕他，只因怕贾母不依，才只 得让他三分。}
}
{
    \eA{while the relationships of father, paternal uncles, and brothers, he did not
        however presume to disregard, as these were among the injunctions bequeathed
        by the holy man,}
}
{
}

\Verse{118}
{
    \zA{现今宝钗生怕宝玉教训他，倒没意思，便连忙替贾环掩饰。}
}
{
    \eA{and he felt bound to listen to a few of their precepts. But to the above
        causes must be assigned the fact that, among his brothers, he did no more
        than accomplish the general purport of the principle of human affections;
        bearing in mind no thought whatever that he himself was a human being of the
        male sex, and that it was his duty to be an example to his younger brothers.}
}
{
}

\Verse{119}
{
    \zA{宝玉道： “大正月里，哭什么？}
}
{
    \eA{And this is why Chia Huan and the others entertained no respect for him,
        though in their veneration for dowager lady Chia,}
}
{
}

\Verse{120}
{
    \zA{这里不好，到别处玩去。}
}
{
    \eA{they yielded to him to a certain degree. Pao-ch'ai harboured fears lest,}
}
{
}

\Verse{121}
{
    \zA{你天天念书，倒念糊涂了。}
}
{
    \eA{on this occasion, Pao-yü should call him to book, and put him out of face,}
}
{
}

\Verse{122}
{
    \zA{譬如这 件东西不好，横竖那一件好，就舍了这件取那件，难道你守着这件东西哭会子就好 了不成？}
}
{
    \eA{and she there and then lost no time in taking Chia Huan's part with a view
        to screening him. "In this felicitous first moon what are you blubbering
        for?" Pao-yü inquired, "if this place isn't nice, why then go somewhere else
        to play.}
}
{
}

\Verse{123}
{
    \zA{你原是要取乐儿，倒招的自己烦恼。}
}
{
    \eA{But from reading books, day after day, you've studied so much that you've
        become quite a dunce.}
}
{
}

\Verse{124}
{
    \zA{还不快去呢！”}
}
{
    \eA{If this thing, for instance,}
}
{
}

\Verse{125}
{
    \zA{贾环听了，只得回来。}
}
{
    \eA{isn't good, that must, of course, be good, so then discard this and take up
        that,}
}
{
}

\Verse{126}
{
    \zA{赵姨娘见他这般，因问：“是那里垫了踹窝来了？”}
}
{
    \eA{but is it likely that by sticking to this thing and crying for a while that
        it will become good? You came originally with the idea of reaping some fun,}
}
{
}

\Verse{127}
{
    \zA{贾 环便说：“同宝姐姐玩来着。}
}
{
    \eA{and you've instead provoked yourself to displeasure, and isn't it better
        then that you should be off at once."}
}
{
}

\Verse{128}
{
    \zA{莺儿欺负我，赖我的钱；宝玉哥哥撵了我来了。”}
}
{
    \eA{Chia Huan upon hearing these words could not but come back to his quarters;
        and Mrs. Chao noticing the frame of mind in which he was felt constrained to
        inquire:}
}
{
}

\Verse{129}
{
    \zA{赵 姨娘啐道：“谁叫你上高台盘了？}
}
{
    \eA{"Where is it that you've been looked down upon by being made to fill up a
        hole, and being trodden under foot?"}
}
{
}

\Verse{130}
{
    \zA{下流没脸的东西!那里玩不得？}
}
{
    \eA{"I was playing with cousin Pao-ch'ai," Chia Huan readily replied, "when Ying
        Erh insulted me,}
}
{
}

\Verse{131}
{
    \zA{谁叫你跑了去讨这 没意思？”}
}
{
    \eA{and deprived me of my money, and brother Pao-yü drove me away." "Ts'ui!"
        exclaimed Mrs. Chao,}
}
{
}

\Verse{132}
{
    \zA{正说着，可巧凤姐在窗外过，都听到耳内，便隔着窗户说道：“大正月 里，怎么了？}
}
{
    \eA{"who bade you (presume so high) as to get up into that lofty tray? You low
        and barefaced thing! What place is there that you can't go to and play; and
        who told you to run over there and bring upon yourself all this shame?"}
}
{
}

\Verse{133}
{
    \zA{兄弟们小孩子家，一半点儿错了，你只教导他，说这样话做什么？}
}
{
    \eA{As she spoke, lady Feng was, by a strange coincidence, passing outside under
        the window; so that every word reached her ear, and she speedily asked from
        outside the window:}
}
{
}

\Verse{134}
{
    \zA{凭他 怎么着，还有老爷太太管他呢，就大口家啐他？}
}
{
    \eA{"What are you up to in this happy first moon? These brothers are, really,
        but mere children, and will you just for a slight mistake,}
}
{
}

\Verse{135}
{
    \zA{他现是主子，不好，横竖有教导他 的人，与你什么相干？}
}
{
    \eA{go on preaching to him! what's the use of coming out with all you've said?
        Let him go wherever he pleases; for there are still our lady and Mr. Chia
        Cheng to keep him in order.}
}
{
}

\Verse{136}
{
    \zA{环兄弟，出来!跟我玩去。”}
}
{
    \eA{But you go and sputter him with your gigantic mouth; he's at present a
        master,}
}
{
}

\Verse{137}
{
    \zA{贾环素日怕凤姐比怕王夫人更甚， 听见叫他，便赶忙出来。}
}
{
    \eA{and if there be anything wrong about him, there are, after all, those to
        rate him; and what business is that of yours? Brother Huan, come out with
        you,}
}
{
}

\Verse{138}
{
    \zA{赵姨娘也不敢出声。}
}
{
    \eA{and follow me and let us go and enjoy ourselves."}
}
{
}

\Verse{139}
{
    \zA{凤姐向贾环道：“你也是个没性气的 东西呦!时常说给你：要吃，要喝，要玩，你爱和那个姐姐妹妹哥哥嫂子玩，就和 那个玩。}
}
{
    \eA{Chia Huan had ever been in greater fear and trembling of lady Feng, than of
        madame Wang, so that when her summons reached his ear, he hurriedly went
        out, while Mrs. Chao, on the other hand, did not venture to breathe a single
        word. "You too," resumed lady Feng, addressing Chia Huan; "are a thing
        devoid of all natural spirit!}
}
{
}

\Verse{140}
{
    \zA{你总不听我的话，倒叫这些人教的你歪心邪意、狐媚魇道的。}
}
{
    \eA{I've often told you that if you want to eat, drink, play, or laugh, you were
        quite free to go and play with whatever female cousin, male cousin, or
        sister-in-law you choose to disport yourself with;}
}
{
}

\Verse{141}
{
    \zA{自己又不尊 重，要往下流里走，安着坏心，还只怨人家偏心呢。}
}
{
    \eA{but you won't listen to my words. On the contrary, you let all these persons
        teach you to be depraved in your heart, perverse in your mind, to be sly,
        artful, and domineering;}
}
{
}
\ChapterImage{img/chapters/068第20-21回 榮國府寶釵做生辰 聽曲文寶玉悟禪機.jpg}


\Verse{142}
{
    \zA{输了几个钱，就这么个样儿！”}
}
{
    \eA{and you've, besides, no respect for your own self, but will go with that
        low-bred lot!}
}
{
}

\Verse{143}
{
    \zA{因问贾环：“你输了多少钱？”}
}
{
    \eA{and your perverse purpose is to begrudge people's preferences! But what
        you've lost are simply a few cash,}
}
{
}

\Verse{144}
{
    \zA{贾环见问，只得诺诺的说道：“输了一二百钱。”}
}
{
    \eA{and do you behave in this manner? How much did you lose?" she proceeded to
        ask Chia Huan; and Chia Huan, upon hearing this question,}
}
{
}

\Verse{145}
{
    \zA{凤姐啐道：“亏了你还是个爷，输了一二百钱就这么着！”}
}
{
    \eA{felt constrained to obey, by saying something in the way of a reply. "I've
        lost," he explained, "some hundred or two hundred cash." "You have,"
        rejoined lady Feng,}
}
{
}

\Verse{146}
{
    \zA{回头叫：“丰儿，去取 一吊钱来；姑娘们都在后头玩呢，把他送了去。}
}
{
    \eA{"the good fortune of being a gentleman, and do you make such a fuss for the
        loss of a hundred or two hundred cash!" and turning her head round, "Feng
        Erh," she added, "go and fetch a thousand cash;}
}
{
}

\Verse{147}
{
    \zA{你明儿再这么狐媚子，我先打了你， 再叫人告诉学里，皮不揭了你的!为你这不尊贵，你哥哥恨得牙痒痒，不是我拦着， 窝心脚把你的肠子还窝出来呢！”}
}
{
    \eA{and as the girls are all playing at the back, take him along to go and play.
        And if again by and by, you're so mean and deceitful, I shall, first of all,
        beat you, and then tell some one to report it at school, and won't your skin
        be flayed for you?}
}
{
}

\Verse{148}
{
    \zA{喝令：“去罢！”}
}
{
    \eA{All because of this want of respect of yours,}
}
{
}

\Verse{149}
{
    \zA{贾环诺诺的，跟了丰儿得了钱， 自去和迎春等玩去，不在话下。}
}
{
    \eA{your elder cousin is so angry with you that his teeth itch; and were it not
        that I prevent him, he would hit you with his foot in the stomach and kick
        all your intestines out!}
}
{
}

\Verse{150}
{
    \zA{且说宝玉正和宝钗玩笑，忽见人说：“史大姑娘来了。”}
}
{
    \eA{Get away," she then cried; whereupon Chia Huan obediently followed Feng Erh,
        and taking the money he went all by himself to play with Ying Ch'un and the
        rest;}
}
{
}

\Verse{151}
{
    \zA{宝玉听了，连忙就走。}
}
{
    \eA{where we shall leave him without another word. But to return to Pao-yü.}
}
{
}

\Verse{152}
{
    \zA{宝钗笑道：“等着，咱们两个一齐儿走，瞧瞧他去。”}
}
{
    \eA{He was just amusing himself and laughing with Pao-ch'ai, when at an
        unexpected moment, he heard some one announce that Miss Shih had come.}
}
{
}

\Verse{153}
{
    \zA{说着，下了炕，和宝玉来至 贾母这边。}
}
{
    \eA{At these words, Pao-yü rose, and was at once going off when "Wait," shouted
        Pao-ch'ai with a smile,}
}
{
}

\Verse{154}
{
    \zA{只见史湘云大说大笑的，见了他两个，忙站起来问好。}
}
{
    \eA{"and we'll go over together and see her." Saying this, she descended from
        the stove-couch, and came, in company with Pao-yü,}
}
{
}

\Verse{155}
{
    \zA{正值黛玉在旁， 因问宝玉：“打那里来？”}
}
{
    \eA{to dowager lady Chia's on this side, where they saw Shih Hsiang-yün laughing
        aloud, and talking immoderately;}
}
{
}

\Verse{156}
{
    \zA{宝玉便说：“打宝姐姐那里来。”}
}
{
    \eA{and upon catching sight of them both, she promptly inquired after their
        healths,}
}
{
}

\Verse{157}
{
    \zA{黛玉冷笑道：“我说 呢!亏了绊住，不然，早就飞了来了。”}
}
{
    \eA{and exchanged salutations. Lin Tai-yü just happened to be standing by, and
        having set the question to Pao-yü "Where do you come from?"}
}
{
}

\Verse{158}
{
    \zA{宝玉道：“只许和你玩，替你解闷儿；不 过偶然到他那里，就说这些闲话。”}
}
{
    \eA{"I come from cousin Pao-ch'ai's rooms," Pao-yü readily replied. Tai-yü gave
        a sardonic smile. "What I maintain is this," she rejoined, "that lucky
        enough for you,}
}
{
}

\Verse{159}
{
    \zA{黛玉道：“好没意思的话!去不去，管我什么 事？}
}
{
    \eA{you were detained over there; otherwise, you would long ago have, at once,
        come flying in here!" "Am I only free to play with you?"}
}
{
}

\Verse{160}
{
    \zA{又没叫你替我解闷儿!还许你从此不理我呢！”}
}
{
    \eA{Pao-yü inquired, "and to dispel your ennui! I simply went over to her place
        for a run, and that quite casually,}
}
{
}

\Verse{161}
{
    \zA{说着，便赌气回房去了。}
}
{
    \eA{and will you insinuate all these things?" "Your words are quite devoid of
        sense,"}
}
{
}

\Verse{162}
{
    \zA{宝玉忙跟了来，问道：“好好儿的又生气了!就是我说错了，你到底也还坐坐 儿，合别人说笑一会子啊？”}
}
{
    \eA{Tai-yü added; "whether you go or not what's that to me? neither did I tell
        you to give me any distraction; you're quite at liberty from this time forth
        not to pay any notice to me!" Saying this, she flew into a high dudgeon and
        rushed back into her room; but Pao-yü promptly followed in her footsteps:}
}
{
}

\Verse{163}
{
    \zA{黛玉道：“你管我呢！”}
}
{
    \eA{"Here you are again in a huff," he urged, "and all for no reason!}
}
{
}

\Verse{164}
{
    \zA{宝玉笑道：“我自然不敢管 你，只是你自己遭塌坏了身子呢。”}
}
{
    \eA{Had I even passed any remark that I shouldn't, you should anyhow have still
        sat in there, and chatted and laughed with the others for a while; instead
        of that, you come again to sit and mope all alone!"}
}
{
}

\Verse{165}
{
    \zA{黛玉道：“我作践了我的身子，我死我的，与 你何干？”}
}
{
    \eA{"Are you my keeper?" Tai-yü expostulated. "I couldn't, of course," Pao-yü
        smiled, "presume to exercise any influence over you; but the only thing is
        that you are doing your own health harm!"}
}
{
}

\Verse{166}
{
    \zA{宝玉道：“何苦来？}
}
{
    \eA{"If I do ruin my health," Tai-yü rejoined, "and I die,}
}
{
}

\Verse{167}
{
    \zA{大正月里，‘死’了‘活’了的。”}
}
{
    \eA{it's my own lookout! what's that to do with you?" "What's the good,"
        protested Pao-yü,}
}
{
}

\Verse{168}
{
    \zA{黛玉道：“偏 说‘死’!我这会子就死!你怕死，你长命百岁的活着，好不好？”}
}
{
    \eA{"of talking in this happy first moon of dying and of living?" "I \_will\_
        say die," insisted Tai-yü, "die now, at this very moment! but you're afraid
        of death; and you may live a long life of a hundred years,}
}
{
}

\Verse{169}
{
    \zA{宝玉笑道：“要 像只管这么闹，我还怕死吗？}
}
{
    \eA{but what good will that be!" "If all we do is to go on nagging in this way,"
        Pao-yü remarked smiling, "will I any more be afraid to die?}
}
{
}

\Verse{170}
{
    \zA{倒不如死了干净。”}
}
{
    \eA{on the contrary, it would be better to die, and be free!" "Quite so!"}
}
{
}

\Verse{171}
{
    \zA{黛玉忙道：“正是了，要是这样 闹，不如死了干净！”}
}
{
    \eA{continued Tai-yü with alacrity, "if we go on nagging in this way, it would
        be better for me to die, and that you should be free of me!" "I speak of my
        own self dying,"}
}
{
}

\Verse{172}
{
    \zA{宝玉道：“我说自家死了干净，别错听了话，又赖人。”}
}
{
    \eA{Pao-yü added, "so don't misunderstand my words and accuse people wrongly."
        While he was as yet speaking, Pao-ch'ai entered the room: "Cousin Shih is
        waiting for you;"}
}
{
}

\Verse{173}
{
    \zA{正 说着，宝钗走来，说：“史大妹妹等你呢。”}
}
{
    \eA{she said; and with these words, she hastily pushed Pao-yü on, and they
        walked away. Tai-yü, meanwhile, became more and more a prey to resentment;}
}
{
}

\Verse{174}
{
    \zA{说着，便拉宝玉走了。}
}
{
    \eA{and disconsolate as she felt, she shed tears in front of the window.}
}
{
}

\Verse{175}
{
    \zA{这黛玉越发气 闷，只向窗前流泪。}
}
{
    \eA{But not time enough had transpired to allow two cups of tea to be drunk,
        before Pao-yü came back again.}
}
{
}

\Verse{176}
{
    \zA{没两盏茶时，宝玉仍来了。}
}
{
    \eA{At the sight of him, Tai-yü sobbed still more fervently and incessantly,}
}
{
}

\Verse{177}
{
    \zA{黛玉见了，越发抽抽搭搭的哭个不住。}
}
{
    \eA{and Pao-yü realising the state she was in, and knowing well enough how
        arduous a task it would be to bring her round,}
}
{
}

\Verse{178}
{
    \zA{宝玉见了这 样，知难挽回，打叠起百样的款语温言来劝慰。}
}
{
    \eA{began to join together a hundred, yea a thousand kinds of soft phrases and
        tender words to console her. But at an unforeseen moment, and before he
        could himself open his mouth,}
}
{
}

\Verse{179}
{
    \zA{不料自己没张口，只听黛玉先说道： “你又来作什么？}
}
{
    \eA{he heard Tai-yü anticipate him. "What have you come back again for?" she
        asked. "Let me die or live, as I please, and have done!}
}
{
}

\Verse{180}
{
    \zA{死活凭我去罢了!横竖如今有人和你玩，比我又会念，又会作，又 会写，又会说会笑，又怕你生气，拉了你去哄着你。}
}
{
    \eA{You've really got at present some one to play with you, one who, compared
        with me, is able to read and able to compose, able to write, to speak, as
        well as to joke, one too who for fear lest you should have ruffled your
        temper dragged you away: and what do you return here for now?"}
}
{
}

\Verse{181}
{
    \zA{你又来作什么呢？”}
}
{
    \eA{Pao-yü, after listening to all she had to say, hastened to come up to her.}
}
{
}

\Verse{182}
{
    \zA{宝玉听了， 忙上前悄悄的说道：“你这么个明白人，难道连‘亲不隔疏，后不僭先’也不知道？}
}
{
    \eA{"Is it likely," he observed in a low tone of voice, "that an intelligent
        person like you isn't so much as aware that near relatives can't be
        separated by a distant relative, and a remote friend set aside an old
        friend! I'm stupid, there's no gainsaying,}
}
{
}

\Verse{183}
{
    \zA{我虽糊涂，却明白这两句话。}
}
{
    \eA{but I do anyhow understand what these two sentiments imply. You and I are,
        in the first place,}
}
{
}

\Verse{184}
{
    \zA{头一件，咱们是姑舅姐妹，宝姐姐是两姨姐妹，论亲 戚也比你远。}
}
{
    \eA{cousins on my father's sister's side; while sister Pao-ch'ai and I are two
        cousins on mother's sides, so that, according to the degrees of
        relationship, she's more distant than yourself.}
}
{
}

\Verse{185}
{
    \zA{第二件，你先来，咱们两个一桌吃，一床睡，从小儿一处长大的，他 是才来的，岂有个为他远你的呢？”}
}
{
    \eA{In the second place, you came here first, and we two have our meals at one
        table and sleep in one bed, having ever since our youth grown up together;
        while she has only recently come, and how could I ever distance you on her
        account?"}
}
{
}

\Verse{186}
{
    \zA{黛玉啐道：“我难道叫你远他？}
}
{
    \eA{"Ts'ui!" Tai-yü exclaimed. "Will I forsooth ever make you distance her!}
}
{
}

\Verse{187}
{
    \zA{我成了什么人 了呢？}
}
{
    \eA{who and what kind of person have I become to do such a thing?}
}
{
}

\Verse{188}
{
    \zA{我为的是我的心！”}
}
{
    \eA{What (I said) was prompted by my own motives."}
}
{
}

\Verse{189}
{
    \zA{宝玉道：“我也为的是我的心。}
}
{
    \eA{"I too," Pao-yü urged, "made those remarks prompted by my own heart's
        motives, and do you mean to say that your heart can only read the feelings
        of your own heart, and has no idea whatsoever of my own?"}
}
{
}

\Verse{190}
{
    \zA{你难道就知道你的心， 不知道我的心不成？”}
}
{
    \eA{Tai-yü at these words, lowered her head and said not a word. But after a
        long interval, "You only know," she continued,}
}
{
}

\Verse{191}
{
    \zA{黛玉听了，低头不语，半日说道：“你只怨人行动嗔怪你， 你再不知道你怄的人难受。}
}
{
    \eA{"how to feel bitter against people for their action in censuring you: but
        you don't, after all, know that you yourself provoke people to such a
        degree, that it's hard for them to put up with it! Take for instance the
        weather of to-day as an example.}
}
{
}

\Verse{192}
{
    \zA{就拿今日天气比，分明冷些，怎么你倒脱了青肷披风 呢？”}
}
{
    \eA{It's distinctly very cold, to-day, and yet, how is it that you are so
        contrary as to go and divest yourself of the pelisse with the bluish
        breast-fur overlapping the cloth?"}
}
{
}

\Verse{193}
{
    \zA{宝玉笑道：“何尝没穿？}
}
{
    \eA{"Why say I didn't wear it?" Pao-yü smilingly observed.}
}
{
}

\Verse{194}
{
    \zA{见你一恼，我一暴燥，就脱了。”}
}
{
    \eA{"I did, but seeing you get angry I felt suddenly in such a terrible blaze,
        that I at once took it off!"}
}
{
}

\Verse{195}
{
    \zA{黛玉叹道：“回 来伤了风，又该讹着吵吃的了。”}
}
{
    \eA{Tai-yü heaved a sigh. "You'll by and by catch a cold," she remarked, "and
        then you'll again have to starve, and vociferate for something to eat!"}
}
{
}

\Verse{196}
{
    \zA{二人正说着，只见湘云走来，笑道：“爱哥哥，林姐姐，你们天天一处玩，我 好容易来了，也不理我理儿。”}
}
{
    \eA{While these two were having this colloquy, Hsiang-yün was seen to walk in!
        "You two, Ai cousin and cousin Lin," she ventured jokingly, "are together
        playing every day, and though I've managed to come after ever so much
        trouble,}
}
{
}

\Verse{197}
{
    \zA{黛玉笑道：“偏是咬舌子爱说话，连个‘二’哥哥 也叫不上来，只是‘爱’哥哥‘爱’哥哥的。}
}
{
    \eA{you pay no heed to me at all!" "It's invariably the rule," Tai-yü retorted
        smilingly, "that those who have a defect in their speech will insist upon
        talking; she can't even come out correctly with 'Erh' (secundus) cousin,}
}
{
}

\Verse{198}
{
    \zA{回来赶围棋儿，又该你闹‘么爱三’ 了。”}
}
{
    \eA{and keeps on calling him 'Ai' cousin, 'Ai' cousin! And by and by when you
        play 'Wei Ch'i' you're sure also to shout out yao,}
}
{
}

\Verse{199}
{
    \zA{宝玉笑道：“你学惯了，明儿连你还咬起来呢。”}
}
{
    \eA{ai, (instead of erh), san; (one, two, three)." Pao-yü laughed. "If you
        imitate her," he interposed, "and get into that habit,}
}
{
}

\Verse{200}
{
    \zA{湘云道：“他再不放人一 点儿，专会挑人。}
}
{
    \eA{you'll also begin to bite your tongue when you talk." "She won't make even
        the slightest allowance for any one,"}
}
{
}

\Verse{201}
{
    \zA{就算你比世人好，也不犯见一个打趣一个。}
}
{
    \eA{Hsiang-yün rejoined; "her sole idea being to pick out others' faults. You
        may readily be superior to any mortal being,}
}
{
}

\Verse{202}
{
    \zA{我指出个人来，你敢 挑他，我就服你。”}
}
{
    \eA{but you shouldn't, after all, offend against what's right and make fun of
        every person you come across! But I'll point out some one,}
}
{
}

\Verse{203}
{
    \zA{黛玉便问：“是谁？”}
}
{
    \eA{and if you venture to jeer her, I'll at once submit to you."}
}
{
}

\Verse{204}
{
    \zA{湘云道：“你敢挑宝姐姐的短处，就算 你是个好的。”}
}
{
    \eA{"Who is it?" Tai-yü vehemently inquired. "If you do have the courage,"
        Hsiang-yün answered, "to pick out cousin Pao-ch'ai's faults,}
}
{
}

\Verse{205}
{
    \zA{黛玉听了冷笑道：“我当是谁，原来是他。}
}
{
    \eA{you then may well be held to be first-rate!" Tai-yü after hearing these
        words, gave a sarcastic smile. "I was wondering,"}
}
{
}

\Verse{206}
{
    \zA{我可那里敢挑他呢？”}
}
{
    \eA{she observed, "who it was. Is it indeed she? How could I ever presume to
        pick out hers?"}
}
{
}

\Verse{207}
{
    \zA{宝玉不等说完，忙用话分开。}
}
{
    \eA{Pao-yü allowed her no time to finish, but hastened to say something to
        interrupt the conversation.}
}
{
}

\Verse{208}
{
    \zA{湘云笑道：“这一辈子我自然比不上你。}
}
{
    \eA{"I couldn't, of course, during the whole of this my lifetime," Hsiang-yün
        laughed, "attain your standard!}
}
{
}

\Verse{209}
{
    \zA{我只保佑着 明儿得一个咬舌儿林姐夫，时时刻刻你可听‘爱’呀‘厄’的去!阿弥陀佛，那时 才现在我眼里呢！”}
}
{
    \eA{but my earnest wish is that by and by should be found for you, cousin Lin, a
        husband, who bites his tongue when he speaks, so that you should every
        minute and second listen to 'ai-ya-os!' O-mi-to-fu, won't then your reward
        be manifest to my eyes!"}
}
{
}

\Verse{210}
{
    \zA{说的宝玉一笑，湘云忙回身跑了。}
}
{
    \eA{As she made this remark, they all burst out laughing heartily,}
}
{
}

\Verse{211}
{
    \zA{要知端详，且听下回分解。}
}
{
    \eA{and Hsiang-yün speedily turned herself round and ran away. But reader, do
        you want to know the sequel?}
}
{
}

\Verse{212}
{
    \zA{(本章完)}
}
{
    \eA{Well, then listen to the explanation given in the next chapter.}
}
{
}
